\subsection{正则扩张}

定义 2. --- 域 K 的一个扩张称为正则的，如果它作为 K-代数是正则的。

命题 4. --- 设 A 是域 K 上的一个代数，且为整环，E 是它的分式域。设 L 是 K 的一个扩张；若环 \(L \otimes_{K} A\) 是整环，则 \(L \otimes_{K} E\) 亦然。

如果 \(L \otimes_{K} A\) 是整环，则它可以嵌入到它的分式域 F 中。记 \(u(x) = x \otimes 1\) （对于 \(x \in L\) ），并用 v 表示 E 到 F 中的、延拓 A 到 F 中的单同态 \(y \mapsto 1 \otimes y\) 的 K-同态。由命题 6（V，第 14 页），F 的子域 \(u(L)\) 和 \(v(E)\) 在 K 上线性不相交；因此，\(u * v\) 作为 \(L \otimes_{K} E\) 到 F 中的同态（V，第 12 页）是单射。这表明 \(L \otimes_{K} E\) 是一个整环。

推论. --- 要使 A 是正则 K-代数，必要且充分条件是它的分式域是 K 的正则扩张。

该条件由命题 4 可知是必要的，由命题 3，a) 可知是充分的。

命题 5. --- 域 K 的每一个纯扩张都是正则的。

由前述推论，只需证明每个多项式代数 \(A = K[X_{i}]_{i \in I}\) 都是正则 K-代数。设 L 是 \(K\) 的一个扩张；环 \(L \otimes_{K} A\) 同构于 \(L[X_{i}]_{i \in I}\) （III，第 449 页，注记 2），因此是整环（IV，第 9 页，命题 8）。

命题 6. --- 设 L 是域 K 的一个扩张。若 L 是正则的，则 L 的每个子扩张都是正则的。反之，若 L 的每个有限生成子扩张都是正则的，则 L 是正则的。

第一个断言由命题 3，a) 得出。

设 M 是 K 的一个扩张，并设 U 是 L 的所有有限生成子扩张的集合。对于每个 \(E \in U\) ，环 \(M \otimes_{K} E\) 可以与 \(M \otimes_{K} L\) 的一个子环等同起来，这样我们就得到 \(M \otimes_{K} L\) 的一个递增有向子环族，其并为 \(M \otimes_{K} L\) 。于是第二个断言立即得证。

命题 7. --- 设 L 是域 K 的一个扩张，M 是一个 L-代数（例如 L 的一个扩张）。若 L 在 K 上正则且 M 是正则 L-代数，则 M 作为 K-代数是正则的。

设 E 是 K 的一个扩张；由于 L 在 K 上是正则的， \(E \otimes_{K} L\) 是整环。于是根据命题 2（V，第 141 页），环 \((\mathrm{E} \otimes_{\mathrm{K}} \mathrm{L}) \otimes_{\mathrm{L}} \mathrm{M}\) 是整环，而与它同构的环 \(E \otimes_{K} M\) （II，第 278 页，命题 2）亦然。由此即得结论。

命题 8. --- 设 L 和 M 是域 K 的两个扩张。

\begin{enumerate}
\def\labelenumi{\alph{enumi})}
\item
  如果 \(\mathbf{M}\) 在 \(\mathbf{K}\) 上是正则的，则整环 \(\mathbf{L} \otimes_{\mathbf{K}} \mathbf{M}\) 的分式域是 \(\mathbf{L}\) 的一个正则扩张。
\item
  若 L 和 M 都是 K 的正则扩张，则 L ⊗\_K M 的分式域亦然。
\end{enumerate}

断言 a) 由命题 3，c)（V，第 141 页）及 V，第 141 页的推论得出；断言 b) 由命题 3，b)（V，第 141 页）及 V，第 141 页的推论得出。

